% !TeX root = ../paper.tex
\section{Problem Statement and Objectives}
\label{sec:problem-statement}

\subsection{Motivation}
% TODO: Describe the motivation behind the project
% TODO: Why is multimodal learning important for tropical and infectious disease diagnosis?
% TODO: What are the current challenges in clinical decision support systems?

\subsection{Problem Statement}
% TODO: Clearly define the problem you are addressing
% TODO: What gaps exist in current approaches?
% TODO: What specific challenges do you aim to solve?

\subsection{Research Objectives}
% TODO: List your main objectives clearly
% TODO: What do you aim to achieve with your system?

\subsection{Research Hypothesis}
% TODO: State your research hypotheses
% State your hypotheses clearly
% What do you expect to find or achieve?

We hypothesize that:

\begin{itemize}
    \item \textbf{H1:} Large multimodal models enhanced with domain-specific retrieval mechanisms (RAG) will outperform standalone fine-tuned models in diagnostic accuracy for rare and complex tropical diseases
    
    \item \textbf{H2:} A hybrid approach combining RAG with targeted fine-tuning will achieve superior performance compared to either method alone, particularly in scenarios requiring both broad medical knowledge and disease-specific expertise
    
    \item \textbf{H3:} Models trained on diverse multimodal data sources will demonstrate better generalization across different clinical presentation patterns compared to single-modality approaches
\end{itemize}

% Adjust hypotheses based on your experimental design and expected outcomes
